%% kapa: exam=1 round=29 year=2018 period=1 form=A subject=real_estate subject_ko=부동산학원론 seq=7 number=87
%% source: Q-Net 감정평가사 기출문제 (cst00309) artlSeq=5202808 file=extracted/1008/5202808/2018년 제29회 감정평가사 1차시험 기출문제/2018년 제29회 감정평가사 1차시험 기출문제/1교시 A형.pdf
%% conversion: kordoc 4.24.51 -> tools/gen1.py (auto); review=proofread-sonnet; flags=none
\begin{question}{87}
입지 및 도시공간구조이론에 관한 설명으로 옳지 않은 것은?
\begin{choices}
  \item 최소마찰비용이론은 경제부문의 집적화 이익이 시공간적으로 누적적 인과 과정을 통해 낙후된 지역까지 파급된다고 본다.
  \item 알론소(Alonso)의 입찰지대곡선은 도심으로부터 교외로 이동하면서 거리에 따라 가장 높은 지대를 지불할 수 있는 산업들의 지대곡선을 연결한 선이다.
  \item 해리스(Harris)와 울만(Ullman)의 다핵심이론은 서로 유사한 활동이 집적하려는 특성이 있다고 본다.
  \item 버제스(Burgess)의 동심원이론은 침입, 경쟁, 천이과정을 수반하는 생태학적 논리에 기반하고 있다.
  \item 호이트(Hoyt)의 선형이론은 도시공간의 성장 및 분화가 주요 교통노선을 따라 확대되면서 나타난다고 본다.
\end{choices}
\end{question}
\answer{1}
